\section{可插拔后端架构}

\subsection{\texttt{pmm\_ops\_t}：后端操作表}

\codefile{src/include/kernel/pmm.h}
\begin{lstlisting}[style=cstyle]
typedef struct pmm_ops {
    const char* name;                    // "bitmap" / "buddy"

    void  (*init)(void);                 // 解析 mmap → 建立数据结构
    void* (*alloc_page)(void);           // 分配一个物理页，返回物理地址
    void  (*free_page)(void* page);      // 释放一个物理页

    unsigned (*total_pages)(void);       // 管理的总页数
    unsigned (*free_pages)(void);        // 当前空闲页数
    unsigned (*managed_base)(void);      // 管理区域的物理起始地址

    unsigned (*page_addr)(unsigned idx); // 页索引 → 物理地址
    int (*page_is_used)(unsigned idx);   // 页索引 → 使用状态
} pmm_ops_t;
\end{lstlisting}

\begin{figure}[H]
  \centering
  % figures/ch07/fig02-pmm.tex — auto-extracted from chapter source
\begin{tikzpicture}[
  block/.style={draw, rounded corners, minimum width=3.5cm, minimum height=0.7cm, font=\small},
  arr/.style={-{Stealth[length=3pt]}, thick},
]
  \node[block, fill=gray!10]   (caller) {调用者 (vmm.c, kmalloc.c)};
  \node[block, fill=blue!10, below=0.8cm of caller] (api) {pmm.h 公开 API};
  \node[block, fill=green!10, below=0.8cm of api]   (disp) {pmm.c dispatch 层};
  \node[block, fill=orange!10, below left=1cm and -0.5cm of disp]  (bm) {pmm\_bitmap.c};
  \node[block, fill=cyan!10, below right=1cm and -0.5cm of disp]   (bd) {pmm\_buddy.c};

  \draw[arr] (caller) -- (api);
  \draw[arr] (api) -- (disp);
  \draw[arr] (disp) -- (bm) node[midway, left, font=\scriptsize] {g\_pmm\_ops};
  \draw[arr] (disp) -- (bd);

  \node[right=0.5cm of bm, font=\scriptsize\itshape] {bitmap 后端};
  \node[right=0.5cm of bd, font=\scriptsize\itshape] {buddy 后端};
\end{tikzpicture}

  \caption{PMM 三层可插拔架构}
\end{figure}

\subsection{Dispatch 层：\texttt{pmm.c}}

\codefile{src/kernel/mm/pmm.c}
\begin{lstlisting}[style=cstyle]
static const pmm_ops_t* g_pmm_ops = NULL;

void pmm_register_backend(const pmm_ops_t* ops) {
    g_pmm_ops = ops;
}

void pmm_init(void) {
    if (!g_pmm_ops) g_pmm_ops = pmm_bitmap_get_ops();  // 默认 bitmap
    printk("[pmm] backend: %s\n", g_pmm_ops->name);
    g_pmm_ops->init();
}

void* pmm_alloc_page(void) {
    if (!g_pmm_ops) return NULL;
    return g_pmm_ops->alloc_page();
}

void pmm_free_page(void* page) {
    if (!g_pmm_ops) return;
    g_pmm_ops->free_page(page);
}
// ... 其他函数同理 ...
\end{lstlisting}

\begin{designbox}
这套"ops 函数指针表 + dispatch 薄层"的模式在本项目中反复出现：
\begin{itemize}
  \item PMM: \texttt{pmm\_ops\_t} → \texttt{pmm.c}
  \item Heap: \texttt{heap\_ops\_t} → \texttt{kmalloc.c}
  \item VMA: \texttt{vma\_ops\_t} → \texttt{vma.c}
\end{itemize}
这是 Linux 内核 \texttt{file\_operations} 模式的简化版——
接口固定，实现可替换，调用者零改动。
\end{designbox}

% ============================================================================

